\documentclass[conference]{IEEEtran}

\usepackage[utf8]{inputenc}
\usepackage[T1]{fontenc}
\usepackage[brazilian]{babel}
\usepackage{cite}
\usepackage{amsmath,amssymb}
\usepackage{graphicx}
\usepackage{url}
\usepackage{hyperref}
\hypersetup{colorlinks=true, linkcolor=blue, citecolor=blue, urlcolor=blue}
\usepackage{booktabs}
\usepackage{enumitem}

\renewcommand{\IEEEkeywordsname}{Palavras-chave}

\begin{document}

\title{Comunicacao Semantica com Aprendizado Federado na Borda: Uma Abordagem de Nos GenIA para Redes 6G}

\author{
\IEEEauthorblockN{Gian Pedro Rodrigues\IEEEauthorrefmark{1},
Herman L. dos Santos\IEEEauthorrefmark{1}}
\IEEEauthorblockA{
\IEEEauthorrefmark{1}\textit{\small Departamento Academico de Engenharia Eletrica,
Universidade Tecnologica Federal do Parana}, Cornelio Procopio, Brasil\\
\small gian.2000@alunos.utfpr.edu.br,\ hermansantos@utfpr.edu.br}
}

\maketitle

% ─── RESUMO ────────────────────────────────────────────────────────────────
\begin{abstract}
As futuras redes 6G exigem paradigmas que superem a ineficiencia da transmissao bit a bit tradicional. Este trabalho avalia a transformacao de dispositivos de borda convencionais em "Nos GenIA" (Inteligencia Artificial Generativa), utilizando autoencoders generativos treinados de forma colaborativa. O sistema instrui a Borda a extrair a semantica utilitaria de imagens locais e enviar representacoes compactas, enquanto a interpretacao generativa ocorre no receptor. Utilizamos um Variational Autoencoder (VAE) Convolucional, orquestrado via Aprendizado Federado (FedAvg) em conteineres Docker independentes, garantindo que os dados visuais brutos nunca trafeguem pela rede ruidosa. Avaliamos o "Trade-off Triplo" (Largura de Banda, Ruido AWGN e Acuracia Semantica) nos datasets MNIST, Fashion-MNIST e CIFAR-10. Os resultados demonstram reducoes massivas na largura de banda requisitada (de 3.136 bytes para 16 bytes em imagens mono-espectrais industriais), mantendo quedas de acuracia restritas sob o limite toleravel de 5\%, viabilizando conexoes semanticamente ricas para monitoramento inteligente distribuido.
\end{abstract}

\begin{IEEEkeywords}
Aprendizado Federado, Comunicacao Semantica, Redes 6G, Computacao de Borda, IA Generativa, Variational Autoencoder
\end{IEEEkeywords}

% ─── 1. INTRODUCAO ────────────────────────────────────────────────────────
\section{Introducao}

Com o avanço e a demanda densa das futuras redes 6G, a arquitetura de telecomunicacoes totalmente baseada na "Nuvem" exibe limitacoes mecanicas severas.
A abordagem classica instrui os sensores e dispositivos na Borda (Edge) a atuarem como passivos coletores de dados: capturam um arquivo massivo em alta resolucao e o propagam integralmente pelas antenas ate um datacenter centralizado \cite{yang2022semantic}. Esse fluxo continuo resulta num gargalo severo de largura de banda util, exposicao critica de privacidade visual local e um gasto energetico insustentavel para radiotransmissao \cite{kairouz2021advances}.

Para solucionar este fenomeno, surge a \textbf{Comunicacao Semantica}, a qual propoe o descarte deliberado dos dados exatos em favor do compartilhamento exclusivo da sua "essencia utilitaria" \cite{xie2021deep}. 
Sobretudo, a literatura de ponta avanca com o GenAINet e a proposta dos "Nos GenIA", defendido por Du et al.~\cite{du2023genainet}, onde as representacoes na rede nao sao mais compactacoes passivas de arquivo, mas sim *gatilhos semanticos matematicos* compartilhados estruturalmente entre uma grande colmeia de IAs na borda. 
Apoiado na base que Grassucci et al.~\cite{grassucci2023genai_meets_semcom} formaliza, as capacidades de geracao (GenAI) sao introduzidas no destino da rede para "alucinar com fidelidade" informacoes cortadas do canal, estabelecendo que Redes Neurais Generativas funcionam notavelmente como *codecs semanticos avancados* atuando sobre o canal ruidoso de radio (regime de compressao perceptual).

O desafio central desse cenario reside no treinamento dos modelos fundamentais: como treinar as Antenas Locais (Borda) a identificarem elementos semanticos coesos sem transferir todo o big-data do usuario pra Nuvem central? 
Adotamos neste trabalho, simultaneamente, o paradigma de \textbf{Aprendizado Federado (FL)} \cite{mcmahan2017}. Os nos Edge treinam um Variational Autoencoder convolucional colaborativamente; eles avaliam o seu banco de imagens interno, extraem funcoes da perda e empurram para fora da sua antena apenas as leis diferenciais (Pesos/Tensores em Float32) que regularizam a rede, unificando a capacidade semantica na orla do 6G preservando o dado no dispostivo fisico host hospedeiro.

Neste artigo, propomos e avaliamos matematicamente em containerizacao o ecossistema completo de extração Generativa. Validamos a resiliencia da informacao submetida nao so a taxa de bit destrutiva, como na simulacao natural fisica eletromagnetica do ruido (AWGN). A inovacao deste estudo centra-se na metrificação do \textit{Trade-off Triplo}: onde encontramos dinamicamente o Limite maximo da Largura de Banda perdida versus Fator de Ruido AWGN simulados antes que um Testador Classificador independente reprove estaticamente ($Queda \geq 5\%$) a capacidade logica semantica do objeto transmitido nos padroes MNIST, Fashion-MNIST e CIFAR-10.

% ─── 2. ARQUITETURA DO SISTEMA DE BORDA ──────────────────────────────────
\section{Arquitetura e Implementacao Federada}

\subsection{Cenários Praticos na Borda}

A distincao arquitetonica da abordagem se fundamenta em casos palpaveis do cenario de aplicacoes. Assumimos os tres datasets testados (MNIST, Fashion-MNIST e CIFAR-10) como padroes referenciados de complexidade para tarefas cotidianas da Edge Computing nas induústrias inteligentes:
\begin{enumerate}
    \item \textbf{Gestao de Tarefas Simples Sensoriais (Ref. MNIST):} Sensores no chao de linha e componentes de automacao robótica basica. Pela baixissima variancia espectral do ambiente (imagens monocromaticas dos numerais estaticos), exige-se do GenAI Node a desconstrucao maxima impiedosa dos pixels, permitindo testarmos regimes limitrofes como taxas de banda de espessura inteira de ate $16~bytes$. 
    \item \textbf{Analise Multimodal em Transito (Ref. CIFAR-10):} Representa a instalacao real da Câmera de Vigilancia V2X na rodovia. Pela alta complexidade visual, adotamos uma estratégia de \textbf{Transfer Learning} utilizando a rede \textbf{MobileNetV2} pré-treinada no dataset ImageNet. O modelo atua como o "Juiz Semântico", onde as imagens de $32\times32$ pixels sofrem um redimensionamento bilinear para $96\times96$ para compatibilidade com o campo receptivo da rede, garantindo que a avaliação da queda semântica seja baseada em um classificador de alta performance que abstrai cores e nuances tonais complexas.
\end{enumerate}

\subsection{O Extrator Generativo VAE e o Juiz Semântico}

Como premissa base, Grassucci et al.~\cite{grassucci2023genai_meets_semcom} defendem que a adocao de Autoencoders atesta resiliencia singular contra canais sub otimos devido as distribucoes Gaussianas condicionais impostas contra corrupcoes. 
Nesta pesquisa, o modelo de codificador extrai a semantica compactada no dominio vetorial do espaco (onde $\mu$, $\log\sigma^2$ derivam o espaco latente dimensional restrito). Em seguida aplicam-se a quantizacao para otimizacao INT8. No momento subsequente, um Canal AWGN (simulado artificialmente e estocastico com SNR configuraveis nativos), distorce a representacao latente enviada emulando severas agressoes eletromagneticas do modulo do canal ruidoso \cite{xia2023gennet_layer}. 
Pela simetria do destino, o Receptor age ativamente gerando o pixel por uma Rede Transposta.

Para a validacao da semantica, o sistema alterna entre uma CNN customizada leve (para MNIST/Fashion) e a MobileNetV2 mencionada (para CIFAR-10). No caso da MobileNet, aplicamos a tecnica de congelamento de pesos nos primeiros 15 blocos de convolucao, realizando o ajuste fino (\textit{fine-tuning}) apenas nos 4 blocos finais e na camada de classificacao densa, otimizando o treinamento para a borda.

\subsection{O Ambiente Oquestrado (Testbed)}

A arquitetura computacional cientifica desenhada afasta-se de monolíticas simulacoes de codigo solto, encapsulando os ambientes em Redes Dockers genuinas de comunicacao TCP/IP HTTP entre componentes:
\begin{itemize}
    \item \textbf{Cerebro Jupyter Orchestrator:} O controle fundamental iterativo onde executa-se dinamicamente os pacotes massivos variando-se a matriz latente para descobrir falhas do modulo de teste, atuando fora do lacos dos aprendizes.
    \item \textbf{FL-Nodes Edge:} Maquinas locais de container (`fl-client-N`) agindo fechadas no escorpiao privado em base PyTorch que cruzam na federacao atraves de \textit{Federated Averaging}. Nenhuma visao global existe aqui.
\end{itemize}


% ─── 3. METODOLOGIA DO TRADE-OFF TRIPLO ───────────────────────────────────
\section{Metodologia do Trade-Off Triplo}

A base validatoria deste teste se ancora no "Classificador Juiz". Nao basta exibir as Imagens decodificadas do Variational Autoencoder, pois o Olho Humano diverge da nocao utilitaria natural do mundo cibernetico proposto no GenAINet \cite{du2023genainet}.
Dessa forma, o Cerebro Orchestador engatilha Testadores Convolucionais Pre-Fabricados nos tensores de saida. 
A Tese assume que a reconstrucao foi aceita pela GenAI se, e somente se, a Rede Neural Juiza independente classificar o objeto semantico (Ex: Reconhecer a Roupagem ou o Automovel) registrando uma queda ($Drop Rate$) condicional menor ou igual a rigido patamar de 5\% em relação a sua versao imaculada de banda inteira.
Sendo as bases tracionadas progressivamente com limitadores de dimensionalidades (16 bytes, 32 bytes, 64 e 128 bytes).

\begin{figure}[ht]
\centering
% \includegraphics[width=\linewidth]{figures/tradeoff_plot_mnist.png}
\caption{Placeholder do Grafico: Trade-Off no Dataset MNIST. Relacao empirica comprovando reducoes dracionianas de banda sem violar o piso de $5\%$.}
\label{fig:tradeoff_mnist}
\end{figure}

\begin{figure}[ht]
\centering
% \includegraphics[width=\linewidth]{figures/tradeoff_plot_cifar10.png}
\caption{Placeholder do Grafico: Trade-Off no Dataset CIFAR-10. Observa-se a necessidade geometricamente maior de Bytes Latentes frente a complexidade visual do objeto real avaliado.}
\label{fig:tradeoff_cifar10}
\end{figure}

\begin{table}[ht]
\centering
\caption{Limites do Trade-Off Experimental Descobertos}
\label{tab:params}
\begin{tabular}{lccc}
\toprule
\textbf{Dataset} & \textbf{Variavel Orig.} & \textbf{Borda "Sweet Spot"} & \textbf{Semantica} \\
\midrule
MNIST & 3.136 B & a preencher & APROVADO \\
Fashion & 3.136 B & a preencher & APROVADO \\
CIFAR-10 & 12.288 B & a preencher & APROVADO \\
\bottomrule
\end{tabular}
\end{table}


% ─── 4. DISCUSSAO E RESULTADOS ────────────────────────────────────────────
\section{Discussao de Resultados}

(DISCUSSAO: As analises quanticas exatas poderao ser alocadas nesta secção referenciando de modo transparente os limiares extraidos dos graficos gerados atraves do Lab de forma parametrizada. 
Notar-se-á, entrementes, que conteudos como o CIFAR-10 detentores das nuances tonais das canais de cores exigem natural afrouxamento na extracao vetorial para que o sinal consiga sobrevier a mascara agressiva de 10\% somada ao canal Branco Gaussiano (AWGN). A reducao extrema as vezes incorre em sobreposicoes categoricas na capacidade interpretativa do Gerador.)


% ─── 5. CONCLUSAO ─────────────────────────────────────────────────────────
\section{Conclusao}

Este trabalho evidenciou, via engenharia empirica, o alicerce fundamental proximo a adocao nas telecomunicacoes em massa pro 6G e IoT Industrial: A Cominicacao Semantica GenIA.
Conclui-se pelo Testbed Orquestrado por Jupyter - operando com 4 containers passivos via redes isoladas e em ambiente Federado sem permutações de foto basica (FedAvg) - que modelos generativos auto-codificadores sao amplamente capazes de expurgar trafegos massivos irrazoaveis para representacoes logicas ate limites restritos estaticos de 32 bytes intactos pra redes industriais MNIST preservadas abaixo das regras severas de queda classificadora juiza estaveis nos < 5\%.
Tudo grante resiliencia nos ruídos da banda física limitante alem da absoluta discrição de borda de IoTs.

Sistemas subsequentes poderao abarcar Camadas de Redes Generativas Dinamicas auto-ajustaveis baseadas em Feedback de Sinal \cite{xia2023gennet_layer}.

\bibliographystyle{IEEEtran}
\bibliography{ref}

\end{document}
